\subsubsection{The Precise Form of Congruent Symmetry: Self-Dual Families and the Boundary Zero Principle}\label{subsec:xi-contractive-symmetry-boundary-zero}
This paragraph presents an operator-theoretic mechanism directly aligned with the unitary slice verification protocol of this paper: the object of symmetry is not a geometric ``axis,'' but rather the automorphism of the projective readout system (Definition \ref{def:xi-read-visible-symmetry}).
In the representation of the completion parameter \(u\), this automorphism can be realized by the involution \(u\mapsto 1/\overline u\) and its conjugation on the kernel family; combined with the congruent verification of strict contraction inside the disk, the verifiable zeros of the Fredholm determinant are structurally compressed onto the boundary slice \(|u|=1\).
\par

\begin{definition}[Congruent symmetric operator family and Fredholm determinant]\label{def:xi-contractive-selfdual-family}
Let \(\mathcal{H}\) be a complex Hilbert space and \(u\in\CC^\times\) be the completion parameter.
Let \(T(u)\) be a family of bounded operators satisfying:
\begin{enumerate}
  \item[\textbf{(A)}] \textbf{Analytic controllability:}
  For \(|u|<1\), \(T(u)\) is trace class and \(u\mapsto T(u)\) is analytic in the unit disk;
  \item[\textbf{(B)}] \textbf{Congruent verification:}
  For \(|u|<1\), one has \(\norm{T(u)}<1\);
  \item[\textbf{(C)}] \textbf{Self-dual symmetry:}
  There exists a fixed anti-linear isometric isomorphism \(J:\mathcal{H}\to\mathcal{H}\) (realizing the ``readout dual'') such that for all \(u\neq 0\),
  \[
  T(u)=J\,T\!\left(\frac{1}{\overline u}\right)^{*}J^{-1}.
  \]
\end{enumerate}
On \(|u|<1\), define the Fredholm determinant
\[
\mathcal{D}_T(u):=\det\!\bigl(I-T(u)\bigr),
\]
where the determinant is taken at \(z=1\) as given in Definition \ref{def:fredholm-determinant}.
By (C), if \(|u|>1\) and \(v:=1/\overline u\), then \(T(u)\) is conjugate to \(T(v)^{*}\), hence \(T(u)\) is also trace class and \(\mathcal{D}_T(u)\) is likewise well-defined.
\end{definition}

\begin{theorem}[Boundary zero principle of congruent symmetry]\label{thm:xi-contractive-boundary-zero}
Under the hypotheses of Definition \ref{def:xi-contractive-selfdual-family}, for any \(u\in\CC^\times\) satisfying \(|u|\neq 1\), the operator \(I-T(u)\) is invertible; consequently the Fredholm determinant satisfies \(\mathcal{D}_T(u)\neq 0\).
Therefore, any verifiable zero capturable by \(\det(I-T(u))\) can only occur on the boundary slice \(|u|=1\).
\end{theorem}

\begin{proof}
If \(|u|<1\), then \(\norm{T(u)}<1\) implies that the Neumann series converges in operator norm and yields
\[
(I-T(u))^{-1}=\sum_{n\ge 0}T(u)^{n}.
\]
Hence \(I-T(u)\) is invertible and \(\mathcal{D}_T(u)\neq 0\).
\par
If \(|u|>1\), let \(v:=1/\overline u\), so that \(|v|<1\).
By the preceding step, \(I-T(v)\) is invertible, hence \(I-T(v)^{*}\) is also invertible.
By the self-dual symmetry (C),
\[
I-T(u)=J\bigl(I-T(v)^{*}\bigr)J^{-1},
\]
so \(I-T(u)\) is conjugate to \(I-T(v)^{*}\) and is therefore also invertible, giving \(\mathcal{D}_T(u)\neq 0\). \qedhere
\end{proof}

\begin{corollary}[Critical slice rigidity (completion mapping caliber)]\label{cor:xi-contractive-critical-slice-rigidity}
Suppose there exists an analytic map \(u=u(s)\) satisfying
\[
|u(s)|=1\quad\Longleftrightarrow\quad \Re(s)=\tfrac12.
\]
If for some \(s_0\) one has \(\mathcal{D}_T(u(s_0))=0\), then necessarily \(\Re(s_0)=\tfrac12\).
In particular, for the map \(u_L(s)=L^{2s-1}\) of this paper (Remark \ref{rem:xi-critical-line-unit-circle}), one has \(\mathcal{D}_T(u_L(s))=0\Rightarrow \Re(s)=\tfrac12\).
\end{corollary}

\begin{proof}
By Theorem \ref{thm:xi-contractive-boundary-zero}, if \(\mathcal{D}_T(u(s_0))=0\), then necessarily \(|u(s_0)|=1\). The hypothesis then gives \(\Re(s_0)=\tfrac12\). \qedhere
\end{proof}

\begin{remark}[The object of symmetry: readout automorphism]\label{rem:xi-symmetry-object-autread}
The ``symmetry'' in Theorem \ref{thm:xi-contractive-boundary-zero} is not a geometric exchange of coordinate axes, but rather an automorphism constraint preserving readout invariance (Definition \ref{def:xi-read-visible-symmetry}):
the involution \(u\mapsto 1/\overline u\) together with the anti-linear conjugation \(J(\cdot)J^{-1}\) transfers the invertibility certified inside the disk to the exterior, thereby compressing verifiable singularities onto the fixed-slice boundary.
This meaning is consistent with the semantic version of the \(\mathbf{PW}\) type safety theorem \ref{thm:xi-offset-null-type-safety}: off-slice degrees of freedom are not addressable under the regime that does not externalize the recording axis.
\end{remark}

\paragraph{Phase-lifted formulation of Ramanujan-type symmetry: boundary unitarity and defect vanishing}
In the Ramanujan/RH-sharp spectral bound caliber of this paper, the ``square-root bound'' confines the nontrivial spectrum within the normalized unit disk (e.g., \S\ref{subsec:pom-rh-sharp} and Remark \ref{rem:kernel-graph-rh}).
To directly align this spectral constraint with the ``invertible closure/residual axis externalization,'' one can introduce a stronger boundary clause: on the unitary slice (\(|u|=1\)), the kernel family no longer carries amplitude contraction but assumes a pure phase (isometric/unitary) form.
For \(|u|\neq 1\), define the non-unitary defect operators
\[
\Delta_T(u):=I-T(u)^{*}T(u),\qquad
\widetilde{\Delta}_T(u):=I-T(u)T(u)^{*}.
\]
When \(\norm{T(u)}\le 1\), both \(\Delta_T(u)\) and \(\widetilde{\Delta}_T(u)\) are positive operators, and \(\Delta_T(u)=0=\widetilde{\Delta}_T(u)\) if and only if \(T(u)\) is a unitary operator.
\par

\begin{theorem}[Defect indices and the degrees of freedom of the minimal unitary dilation]\label{thm:xi-defect-residual-degree}
Let \(T\) be a contraction operator (\(\norm{T}\le 1\)) on a Hilbert space \(\mathcal{H}\), and define the defect spaces
\[
\mathcal{D}_{+}(T):=\overline{\mathrm{Ran}}\,(I-T^{*}T)^{1/2},
\qquad
\mathcal{D}_{-}(T):=\overline{\mathrm{Ran}}\,(I-TT^{*})^{1/2}.
\]
Then for any unitary dilation \(U\) acting on \(\mathcal{K}\supseteq \mathcal{H}\) and satisfying
\(
T=P_{\mathcal{H}}U|_{\mathcal{H}}
\), one necessarily has
\[
\dim(\mathcal{K}\ominus\mathcal{H})\ \ge\ \max\{\dim\mathcal{D}_{+}(T),\ \dim\mathcal{D}_{-}(T)\}.
\]
In the finite-rank/finite-dimensional caliber, this lower bound is attainable; in particular, when \(\dim\mathcal{D}_{+}(T)=\dim\mathcal{D}_{-}(T)\), there exists a unitary dilation of minimal dimension such that
\(
\dim(\mathcal{K}\ominus\mathcal{H})=\dim\mathcal{D}_{+}(T)
\).
\end{theorem}

\begin{proof}
This result is a standard fact from the unitary dilation theory of contraction operators (see \cite{SzNagyFoias1970HarmonicAnalysis}); for convenience of alignment with the ``minimal residual axis/minimal externalized degrees of freedom'' certificate semantics of this paper, a direct dimension lower bound computation is provided here.
Let \(V:\mathcal{H}\hookrightarrow\mathcal{K}\) be the natural embedding, so that \(T=V^{*}UV\) and
\[
I-T^{*}T
=V^{*}U^{*}(I-VV^{*})UV
=V^{*}U^{*}P_{\mathcal{K}\ominus\mathcal{H}}UV.
\]
Let \(A:=P_{\mathcal{K}\ominus\mathcal{H}}UV:\mathcal{H}\to\mathcal{K}\ominus\mathcal{H}\), then \(I-T^{*}T=A^{*}A\), whence
\(
\dim\mathcal{D}_{+}(T)\le \dim(\mathcal{K}\ominus\mathcal{H})
\).
Similarly,
\[
I-TT^{*}
=V^{*}U(I-VV^{*})U^{*}V
=V^{*}UP_{\mathcal{K}\ominus\mathcal{H}}U^{*}V
\]
can also be written in the form \(B^{*}B\) for some operator, whence
\(
\dim\mathcal{D}_{-}(T)\le \dim(\mathcal{K}\ominus\mathcal{H})
\).
Combining yields the stated lower bound.
\par
In the finite-dimensional case, one can complete the column (or row) vectors of the matrix \(T\) to an orthonormal set in minimal additional dimensions, thereby constructing a unitary matrix extension attaining the lower bound. \qedhere
\end{proof}

\begin{corollary}[Boundary criterion for defect vanishing]\label{cor:xi-defect-vanishing-boundary}
If the boundary limit at \(|u|=1\) satisfies \(\Delta_T(u)=0\) (equivalently, \(T(u)\) is isometric and surjective on the closure), then there are no additional degrees of freedom that must be supported by continuous amplitude residuals: achieving invertibility at that point does not require externalizing any continuous orthogonal register.
Combined with Theorem \ref{thm:xi-contractive-boundary-zero}: under the congruent symmetry regime, verifiable zeros can only fall on \(|u|=1\), and their boundary dynamics can be supported by pure phase (unitary) degrees of freedom.
\end{corollary}

\paragraph{Structural breakdown away from the fixed slice (falsification/upgrade guardrails)}
Within the above framework, if a class of ``zero-point events'' is claimed to occur at \(|u|\neq 1\), then at least one of the following structures must fail:
\begin{enumerate}
  \item \textbf{Congruence breakdown:} there exists \(|u|<1\) with \(\norm{T(u)}\ge 1\), so that the disk-interior invertibility certificate fails;
  \item \textbf{Self-duality breakdown:} the readout automorphism conjugation of (C) cannot be maintained, so that disk-exterior invertibility cannot be transferred from the disk interior;
  \item \textbf{Interface extension:} additional mode axes are introduced to support radial/amplitude residuals, so that ``zeros'' are no longer captured by the verifiable determinant \(\det(I-T(u))\) but transfer to extended spectral singularities (consistent with the interface dichotomy of Proposition \ref{prop:xi-offcritical-falsifiable-restatement}).
\end{enumerate}
Therefore, ``whether zeros can leave the fixed slice'' can be rewritten under the present regime as a purely structural proposition: whether one can maintain the invertibility of the Fredholm determinant caliber without extending the mode axes and while preserving both the congruent verification and the readout automorphism.
\par

\begin{theorem}[Congruent symmetry--defect vanishing--fixed slice support principle]\label{thm:xi-contract-defect-critical-slice-principle}
Suppose \(T(u)\) satisfies conditions (A)--(C) of Definition \ref{def:xi-contractive-selfdual-family}, and satisfies defect vanishing at the boundary limit \(|u|=1\) (the condition of Corollary \ref{cor:xi-defect-vanishing-boundary}).
Then all verifiable zeros captured by \(\mathcal{D}_T(u)=\det(I-T(u))\) can only lie on the unitary slice \(|u|=1\).
If one further takes a completion map \(u=u(s)\) satisfying \(|u(s)|=1\iff \Re(s)=\tfrac12\), then the corresponding zeros must fall on the critical slice \(\Re(s)=\tfrac12\).
\end{theorem}

\endinput
